\documentclass[12pt,letterpaper]{article}

% --- Paquetes Requeridos ---
\usepackage[utf8]{inputenc}
\usepackage[spanish,es-tabla]{babel}
\usepackage[margin=2.5cm]{geometry}
\usepackage{graphicx}
\usepackage{float}
\usepackage{titlesec}
\usepackage{setspace}
\usepackage{hyperref}
\usepackage{booktabs}
\usepackage{caption}
\usepackage{enumitem}
\usepackage{listings}
\usepackage{xcolor}

% --- Configuración de Párrafo ---
\onehalfspacing
\setlength{\parindent}{0pt}
\setlength{\parskip}{1em}

% --- Rutas de búsqueda de imágenes ---
% Permite compilar tanto desde Documentos/Entregables/ como con todas las
% imágenes en una sola carpeta plana.
\graphicspath{{./}{../Evidencia/}{Evidencia/}{imagenes/}}

% --- Configuración de listings ---
\lstset{
    basicstyle=\ttfamily\footnotesize,
    breaklines=true,
    frame=single,
    captionpos=b,
    showstringspaces=false
}

% --- Inicio del Documento ---
\begin{document}

% --- PORTADA ---
\begin{titlepage}
    \centering

    {\Large \textbf{UNIVERSIDAD SANTO TOMÁS}} \\[1cm]

    {\Large \textbf{Sistema colaborativo de robots móviles para asistencia de orientación en entornos interiores con múltiples pisos}} \\[2cm]

    {\large \textbf{Realizado por:}} \\
    {\large Santiago Hernández Ávila} \\
    {\large Jonny Alejandro Mejía León} \\[1cm]

    \begin{figure}[H]
        \centering
        \includegraphics[width=0.6\textwidth]{Logos Uni.png}
        \label{Logo Universidad Santo Tomas y facultad de ingenieria electronica}
    \end{figure}

    \vfill

    {\large \textbf{Informe de Avance --- Semana 27}} \\
    {\large Fase 6: campaña, análisis y escritura} \\
    {\large Protocolo completo y campaña con los dos vehículos} \\[1.5cm]

    \textbf{Dirigido por:} \\
    Ing. Armando Mateus Rojas, Msc. \\
    Ing. Nestor Ivan Ospina, Msc. \\
    Ing. Oscar Mauricio Gélvez Lizarazo, Msc. \\[1.0cm]

    \vfill

    {\large GED (Grupo de Estudio y Desarrollo en Robótica)} \\
    {\large Facultad de Ingeniería Electrónica} \\
    {\large División de Ingenierías} \\
    {\large Bogotá, octubre de 2026}
\end{titlepage}

% --- CONTENIDO ---

\section{Introducción}

Este informe presenta el avance de la semana 27 del cronograma (12 al 18 de octubre de 2026), la última de toma de datos. El corte C-3 del acta de la decisión GO/NO-GO, el viernes 16 de octubre, pedía dos compuertas: G-5, una misión con relevo entre los dos vehículos, y G-6, las repeticiones de RF-27, el requisito de la demostración física con entre 5 y 10 misiones. Una misión consiste en guiar a un usuario de un origen a un destino; el relevo es el traspaso del usuario del vehículo de un piso al vehículo del otro, en la escalera~\cite{acta}.

Las dos compuertas se alcanzaron el 16 de octubre en los pisos 3 y 4 del edificio. Nueve misiones del sorteo corrieron y se completaron con relevo, y la décima no terminó porque el computador portátil del coordinador se apagó por batería. Con RF-27 verificado, los cuatro objetivos específicos tienen todos sus requisitos verificados.

\begin{table}[H]
    \centering
    \small
    \caption{Compuertas del acta GO/NO-GO al 18 de octubre.}
    \label{tab:compuertas}
    \begin{tabular}{@{}p{4.2cm}p{3.6cm}p{6.2cm}@{}}
        \toprule
        \textbf{Compuerta} & \textbf{Qué exige} & \textbf{Estado al 18 de octubre} \\ \midrule
        G-1, actuación & El vehículo responde a las órdenes de dirección y tracción & Alcanzada el 22 de septiembre \\ \addlinespace
        G-2, odometría & Error de la odometría de hasta 10~\% sobre un recorrido medido de al menos 5~m & Alcanzada el 2 de octubre \\ \addlinespace
        G-3, navegación de un vehículo & Un recorrido de punto a punto con la llegada verificada & Alcanzada el 7 de octubre \\ \addlinespace
        G-4, coexistencia & Los dos vehículos y el coordinador en la misma red, sin conflicto de nombres & Alcanzada el 22 de septiembre \\ \addlinespace
        G-5, protocolo completo & Una misión con relevo entre los dos vehículos, con su registro & Alcanzada el 16 de octubre \\ \addlinespace
        G-6, RF-27 & Entre 5 y 10 misiones del protocolo completo, con su registro & Alcanzada el 16 de octubre, con 9 misiones \\ \bottomrule
    \end{tabular}
\end{table}

\section{Objetivos de la semana}

El cronograma fija para esta semana el siguiente criterio de cierre~\cite{cronograma,plan}:

\begin{quote}
\textit{Presentación completa y ensayada al menos una vez de principio a fin.}
\end{quote}

El criterio no se cumplió: la presentación no se armó en la semana. El 25 de septiembre, al replanificar contra las fechas del acta, la campaña de RF-27 pasó a esta semana, y el trabajo se organizó en \texttt{PLAN\_S27.md} alrededor del corte C-3. La otra actividad del cronograma, las conclusiones sobre la viabilidad técnica, tiene ya los datos que la sostienen (secciones 5 y 6).

\section{Preparación de la campaña}

Las diez misiones de G-6 se sortearon el 13 de octubre, antes de cualquier corrida, como pide el protocolo experimental. Se usó la semilla 20261013, con cinco misiones del piso 3 al 4 y cinco del 4 al 3 entre los salones y el ascensor de cada piso, sin repetir pares y con no más de tres seguidas en el mismo sentido (\texttt{campana\_rf27\_misiones.csv})~\cite{plan,estado}.

El mismo día se corrigió la verificación de llegada del coordinador. Con las grabaciones del 7 de octubre se comprobó que, cuando el sistema de navegación (Nav2) daba la meta por alcanzada, el vehículo seguía moviéndose y la pose (posición y orientación) leída en ese instante lo ubicaba entre 0,53~m y 0,96~m de la meta; un segundo después estaba entre 0,09~m y 0,53~m. El coordinador espera ahora 1,5~s antes de comparar la pose con el punto de destino. La prueba del coordinador con vehículos simulados pasa 19 de 19 escenarios, incluidos el de la inercia y el de la media vuelta que el coordinador pide al vehículo cuando el destino queda detrás.

\section{Revisión de los requisitos de la plataforma robótica}

El 15 de octubre se revisaron los requisitos RF-11 a RF-14 y RF-16 del segundo objetivo específico (OE2, la plataforma robótica con dos vehículos) contra la prueba escrita de cada uno y la evidencia de campo ya registrada: G-2, G-3 y la prueba del sistema completo en el laboratorio del 9 de octubre. Los cinco pasaron a verificados. OE2 quedó con sus seis requisitos verificados~\cite{requisitos}.

\section{Prueba del protocolo completo (G-5)}

La primera sesión del 16 de octubre, entre las 17:38 y las 17:56, corrió doce misiones elegidas por el equipo, sin seguir el sorteo. Se trata como una sesión exploratoria. De ella sale G-5: la misión del Salón 302, en el piso 3, al Salón 402, en el piso 4, se completó con relevo en 132,5~s, con un error de llegada de 0,444~m y 61,1~s entre el final del primer tramo y el inicio del segundo. El registro de la misión es \texttt{G5\_01.json}~\cite{campana}.

\section{Campaña de RF-27 (G-6)}

La segunda sesión, entre las 18:30 y las 19:35, corrió las misiones del sorteo en su orden, sin reubicar los vehículos entre una y otra. El vehículo \texttt{amss-ez9n} atendió el piso 3 y \texttt{amss-jgm9} el piso 4; el coordinador y la grabación corrieron en el portátil. Las misiones 1 a 9 se completaron, todas con relevo~\cite{campana,metricas}.

\begin{table}[H]
    \centering
    \small
    \caption{Las nueve misiones de la campaña.}
    \label{tab:misiones}
    \begin{tabular}{@{}p{0.9cm}p{4.3cm}p{2.0cm}p{2.1cm}p{2.3cm}p{2.3cm}@{}}
        \toprule
        \textbf{N.\textsuperscript{o}} & \textbf{Origen $\rightarrow$ destino} & \textbf{Tiempo total (s)} & \textbf{Respuesta (s)} & \textbf{Espera en el relevo (s)} & \textbf{Error de llegada (m)} \\ \midrule
        1 & Ascensor (piso 3) $\rightarrow$ Salón 403 & 52,3 & 0,51 & 21,06 & 0,435 \\ \addlinespace
        2 & Salón 302 $\rightarrow$ Salón 401 & 90,1 & 0,27 & 0,85 & 0,202 \\ \addlinespace
        3 & Salón 401 $\rightarrow$ Salón 302 & 87,6 & 3,90 & 12,92 & 0,143 \\ \addlinespace
        4 & Salón 403 $\rightarrow$ Salón 303 & 78,9 & 1,25 & 0,71 & 0,420 \\ \addlinespace
        5 & Salón 402 $\rightarrow$ Salón 301 & 105,4 & 0,66 & 0,68 & 0,285 \\ \addlinespace
        6 & Salón 301 $\rightarrow$ Salón 401 & 89,8 & 2,97 & 27,48 & 0,226 \\ \addlinespace
        7 & Ascensor (piso 4) $\rightarrow$ Salón 301 & 81,7 & 0,87 & 1,22 & 0,234 \\ \addlinespace
        8 & Salón 302 $\rightarrow$ Salón 403 & 81,8 & 0,95 & 26,76 & 0,214 \\ \addlinespace
        9 & Salón 303 $\rightarrow$ Salón 402 & 92,6 & 0,63 & 0,73 & 0,299 \\ \bottomrule
    \end{tabular}
\end{table}

El tiempo de respuesta va de la solicitud al primer movimiento del vehículo. La espera en el relevo va desde que el vehículo del primer piso termina su tramo hasta que el del segundo empieza el suyo, y depende de cuánto tarda la persona en subir o bajar y confirmar.

\begin{table}[H]
    \centering
    \small
    \caption{Métricas de la campaña.}
    \label{tab:metricas}
    \begin{tabular}{@{}p{6.0cm}p{8.0cm}@{}}
        \toprule
        \textbf{Métrica} & \textbf{Resultado} \\ \midrule
        Tasa de éxito (RF-23) & 9 de 9, intervalo de confianza del 95~\% (Wilson) del 70,1~\% al 100~\% \\ \addlinespace
        Tiempo de respuesta (RF-21) & Mediana 0,87~s, de 0,27~s a 3,90~s \\ \addlinespace
        Tiempo de asignación (RF-22) & 0,3~ms, con el reloj del portátil \\ \addlinespace
        Continuidad entre niveles (RF-24) & 9 de 9 misiones sin interrupción del guiado; espera mediana de 1,22~s \\ \addlinespace
        Error de llegada & Mediana 0,234~m, de 0,143~m a 0,435~m; las nueve por debajo de 0,5~m \\ \addlinespace
        Tiempo total de la misión & Mediana 87,6~s, de 52,3~s a 105,4~s \\ \bottomrule
    \end{tabular}
\end{table}

Con nueve misiones, lo que el intervalo permite afirmar es que la tasa de éxito supera el 70~\%.

Se declaran tres desviaciones del protocolo, decididas por el equipo el 18 de octubre:

\begin{enumerate}[leftmargin=*]
    \item El error de llegada es el de la localización del propio vehículo (AMCL, la localización sobre el mapa con el LiDAR, el sensor láser de distancias), comparada con el punto de destino 1,5~s después de que Nav2 diera la meta por alcanzada. El protocolo pide una medida externa con flexómetro sobre marcas en el piso, y en esta campaña no se tomó. La referencia disponible es la del 7 de octubre en el piso 4~\cite{piso4}, donde AMCL quedó a 0,10~m o menos de la medida con flexómetro en el avance y a 5 cm o menos en lo lateral. Los registros llevan la fuente de la medida, y el análisis la declara.
    \item Las misiones 1 y 2 se completaron en un segundo intento. En el primero, la 1 terminó con el vehículo a 0,526~m del Ascensor, fuera de la tolerancia de 0,5~m, y la 2 sin la pose necesaria para verificar la llegada. Esos intentos no entran en la tasa.
    \item La campaña se grabó en un solo archivo, partido después por misión con la ventana de tiempo de cada registro del coordinador.
\end{enumerate}

La misión 10 (Salón 403 $\rightarrow$ Ascensor del piso 3) no se cuenta: el portátil se apagó por batería 18~s después de pedirla. En un intento anterior se eligió por error el Ascensor del piso 4 como destino; ese intento no es una misión del sorteo.

\section{Incidencias de la campaña}

Antes de la primera misión válida, entre las 18:34 y las 18:58, ocho intentos de la misión 1 fallaron con la navegación detenida en los dos vehículos. En \texttt{amss-ez9n} el navegador tenía detenida la transformada entre el mapa y la odometría (la relación entre los dos marcos de referencia) desde las 18:21, aunque la localización la seguía publicando; en \texttt{amss-jgm9} el gestor de Nav2 lo apagó porque el planificador dejó de responder durante 60~s. Se reiniciaron las cadenas de navegación de los dos vehículos, y la de \texttt{amss-jgm9} una segunda vez. Desde las 19:10 las nueve misiones corrieron sin intervenir los vehículos. La causa de las transformadas detenidas no se estableció; coincidió con cadenas que llevaban horas encendidas y con la carga de las tarjetas entre 10 y 14~\cite{campana}.

El portátil se apagó dos veces por batería, a las 17:58 y a las 19:35. Las dos grabaciones se recuperaron, y las originales quedan respaldadas con sus sumas de verificación.

\section{Cambios en la interfaz y en las herramientas de registro}

En la etapa de espera del relevo, la página mostraba ``¿Ya subió?'' también cuando el usuario bajaba del piso 4 al 3. Ahora el título depende del sentido (``Suba con calma al piso 4'', ``Baje con calma al piso 3'') y el botón nombra el piso. El título de una misión fallida pasó de ``Camino bloqueado'' a ``No se pudo completar'', porque la misión también falla cuando el vehículo se detiene fuera de la tolerancia. En campo, sin embargo, el teléfono siguió mostrando la versión anterior: el servidor de la página solo enviaba la fecha de cada archivo, y con eso el navegador usó la copia que tenía guardada sin consultar al servidor, incluso al recargar. Se reprodujo el 18 de octubre en el navegador del portátil y se corrigió con un servidor que pide al navegador consultar cada archivo antes de usar su copia. Los textos nuevos quedan por verse en el teléfono~\cite{esquema}.

Para la campaña sin flexómetro, el programa que compone los registros admite una segunda fuente de la llegada, la localización del agente, que queda escrita en cada registro y nunca como medida con flexómetro. Una herramienta nueva parte una grabación de varias misiones en una por misión.

\section{Documento final}

Los capítulos 2 (marco referencial) y 3 (desarrollo) se reescribieron. El marco referencial se apoyó en la literatura, con 29 referencias nuevas verificadas, y el desarrollo amplió las etapas y los cuatro objetivos con tres figuras nuevas. El documento se revisó completo para que el lector de la tesis no dependa de los códigos internos del proyecto. Desde el 15 de octubre compila en el equipo del proyecto: 79 páginas, sin errores ni referencias sin resolver~\cite{estado}.

\section{Aporte a los objetivos específicos}

\begin{table}[H]
    \centering
    \small
    \caption{Aporte de las actividades a los objetivos.}
    \label{tab:aporte}
    \begin{tabular}{@{}p{4.2cm}p{3.6cm}p{6.2cm}@{}}
        \toprule
        \textbf{Actividad} & \textbf{Objetivo} & \textbf{Aporte} \\ \midrule
        Revisión de RF-11 a RF-14 y RF-16 & OE2 & La plataforma con dos vehículos queda con sus seis requisitos verificados \\ \addlinespace
        Prueba del protocolo completo & OE2, OE4 & G-5 alcanzada con una misión con relevo entre los dos vehículos \\ \addlinespace
        Campaña de RF-27 & OE4 & Las cuatro métricas medidas en los vehículos, con 9 misiones \\ \addlinespace
        Confirmación del cambio de piso según el sentido & OE3 & La interfaz indica si el usuario sube o baja \\ \bottomrule
    \end{tabular}
\end{table}

\begin{table}[H]
    \centering
    \small
    \caption{Avance por objetivo específico.}
    \label{tab:avance}
    \begin{tabular}{@{}p{2.0cm}p{2.0cm}p{5.2cm}p{4.8cm}@{}}
        \toprule
        \textbf{Objetivo} & \textbf{Avance} & \textbf{Aporte de la semana} & \textbf{Pendiente} \\ \midrule
        OE1 & 100~\% & --- & --- \\ \addlinespace
        OE2 & 100~\% & RF-11 a RF-14 y RF-16 verificados & --- \\ \addlinespace
        OE3 & 100~\% & Texto de la confirmación según el sentido & --- \\ \addlinespace
        OE4 & 100~\% & RF-27 verificado con 9 misiones & --- \\ \bottomrule
    \end{tabular}
\end{table}

El avance de cada objetivo es el promedio de sus requisitos según su estado en la matriz de requisitos. Con RF-27 verificado, el promedio de los cuatro objetivos pasa del 92,0~\% de la semana 26 al 100~\%, frente al 84,4~\% del calendario (semana 27 de 32). De la matriz queda parcial la restricción RNF-03, que pide registrar como observación las personas que pasan durante las pruebas~\cite{requisitos}.

\section{Estado del cronograma y trabajo previsto}

\begin{table}[H]
    \centering
    \small
    \caption{Criterio de cierre de la semana 27.}
    \label{tab:criterio}
    \begin{tabular}{@{}p{4.2cm}p{3.6cm}p{6.2cm}@{}}
        \toprule
        \textbf{Parte del criterio de cierre} & \textbf{Estado} & \textbf{Evidencia} \\ \midrule
        Presentación completa y ensayada & No cumplida & --- \\ \addlinespace
        Conclusiones sobre la viabilidad técnica & Con los datos de la campaña; sin redactar & \texttt{S27\_campana\_rf27.md} \\ \bottomrule
    \end{tabular}
\end{table}

El corte C-3 cerró la toma de datos el 16 de octubre. Para las semanas siguientes el equipo prevé~\cite{cronograma}:

\begin{enumerate}[leftmargin=*]
    \item Redactar en el documento final los resultados del sistema completo y las conclusiones.
    \item Registrar las observaciones de personas en el pasillo durante la campaña, para RNF-03.
\end{enumerate}

\section{Conclusiones}

\begin{enumerate}[leftmargin=*]
    \item G-5 y G-6 se alcanzaron el 16 de octubre: nueve misiones del sorteo se completaron con relevo entre los dos vehículos, cada una con su registro.
    \item Las nueve misiones dan 9 de 9 éxitos, con un intervalo de confianza del 95~\% del 70,1~\% al 100~\%, y un error de llegada mediano de 0,234~m según la localización del vehículo, sin medida externa.
    \item El sistema respondió a cada solicitud en una mediana de 0,87~s y asignó el vehículo en 0,3~ms; una misión entre pisos tomó una mediana de 87,6~s.
    \item Con RF-27 verificado, los cuatro objetivos específicos tienen todos sus requisitos verificados.
    \item Antes de la campaña, la navegación se detuvo en los dos vehículos y hubo que reiniciarla; la causa queda sin establecer.
\end{enumerate}

\section{Anexos}

Las cifras de este informe proceden de los registros y documentos de evidencia del repositorio del proyecto, que se listan a continuación. El mismo informe está disponible en Markdown, en \texttt{Entregable\_semana\_27.md}, para leerlo directamente en el repositorio.

\begin{thebibliography}{99}

\bibitem{acta}
S. Hernández Ávila y J. A. Mejía León,
``Acta de decisión GO/NO-GO de la demostración física,''
Repositorio del proyecto, \texttt{Documentos/ACTA\_GO\_NOGO.md}, octubre de 2026.

\bibitem{plan}
S. Hernández Ávila,
``Plan de la semana 27,''
Repositorio del proyecto, \texttt{Documentos/PLAN\_S27.md}, octubre de 2026.

\bibitem{requisitos}
S. Hernández Ávila y J. A. Mejía León,
``Matriz de requisitos,''
Repositorio del proyecto, \texttt{Documentos/REQUISITOS.md}, octubre de 2026.

\bibitem{campana}
S. Hernández Ávila,
``Campaña del sistema completo con los dos vehículos, pisos 3 y 4,''
Repositorio del proyecto, \texttt{Documentos/Evidencia/S27\_campana\_rf27.md}, octubre de 2026.

\bibitem{metricas}
S. Hernández Ávila,
``Métricas de la campaña de RF-27,''
Repositorio del proyecto, \texttt{Documentos/Evidencia/S27\_metricas\_rf27.json}, octubre de 2026.

\bibitem{piso4}
S. Hernández Ávila,
``Cadena del piso 4 y media vuelta en la salida,''
Repositorio del proyecto, \texttt{Documentos/Evidencia/S26\_piso4\_cadena\_media\_vuelta.md}, octubre de 2026.

\bibitem{esquema}
S. Hernández Ávila,
``Esquema del registro de misión,''
Repositorio del proyecto, \texttt{Documentos/ESQUEMA\_REGISTRO\_MISION.md}, octubre de 2026.

\bibitem{estado}
S. Hernández Ávila y J. A. Mejía León,
``Tablero de estado del proyecto,''
Repositorio del proyecto, \texttt{ESTADO.md}, octubre de 2026.

\bibitem{cronograma}
S. Hernández Ávila y J. A. Mejía León,
``Cronograma actualizado, semanas 17 a 32,''
Repositorio del proyecto, \texttt{Documentos/CRONOGRAMA\_S17\_S32.md}, agosto de 2026.

\end{thebibliography}

\end{document}
